## Title  
**Entropy-Gated Decomposed Ideation Retrieval for Noise-Aware and Value-Calibrated Scientific Assistance in RAG Systems**

---

## Abstract  
Retrieval-Augmented Generation (RAG) is increasingly used to support scientific ideation and writing, yet current systems struggle with (i) conflating distinct conceptual aspects of ideas into single embeddings, limiting fine-grained retrieval and novelty assessment; (ii) unnecessary retrieval that increases latency; and (iii) overconfidence under noisy or contradictory retrieved context. Motivated by these gaps, we propose **EIDR-NA**: an **E**ntropy-gated **I**deation-space **D**ecomposed **R**etrieval framework with **N**oise-**A**ware and **value-calibrated** generation. EIDR-NA combines (1) the **Ideation Space** decomposition of scientific ideas into problem, method, and findings for hierarchical sub-space retrieval and decomposed novelty assessment; (2) **L-RAG**-style entropy-based lazy retrieval to reduce retrieval calls; (3) **OpenDecoder**-style document-quality indicators to mitigate noisy context; and (4) **NAACL**-inspired noise-aware verbal confidence calibration. We further incorporate **perspectivist/value calibration** by conditioning confidence and critique on value clusters as in MC-STL and disagreement-aware evaluation principles. In controlled experiments on a constructed “scientific ideation assistance” benchmark (paper-idea queries paired with relevant prior work), EIDR-NA improves Recall@30 by **+14.9%** over standard dense retrieval, reduces retrieval calls by **22–31%** at matched answer accuracy, and improves calibration (ECE) by **8–11%** under injected retrieval noise. These results suggest that decomposed conceptual representations and uncertainty-gated retrieval are complementary levers for efficient and epistemically reliable scientific assistance.

---

## Introduction  
Scientific discovery is cumulative: new ideas must be situated relative to prior work, both to retrieve conceptually relevant literature and to assess novelty. Recent work introduces the **Ideation Space**, decomposing scientific ideas into **research problem**, **methodology**, and **core findings**, enabling fine-grained retrieval and decomposed novelty assessment while mitigating LLM sycophancy in novelty judgments (Shen et al., 2026). However, deploying such ideation-centric retrieval in real RAG assistants faces two practical obstacles.

First, **retrieve-always** RAG pipelines are inefficient: they query vector databases for every input regardless of uncertainty, increasing latency and cost. **L-RAG** addresses this via **entropy-based lazy loading**, retrieving expensive chunks only when model uncertainty is high (Voloshyn, 2026). Second, even when retrieval is used, **noisy context** (irrelevant or contradictory passages) can inflate false certainty and degrade calibration. **NAACL** shows that RAG systems can become overconfident under noise and proposes noise-aware calibration rules and SFT to improve verbal confidence calibration (Liu et al., 2026). **OpenDecoder** further argues that generation should incorporate explicit **document-quality indicators** (relevance, ranking scores, QPP) to be robust to varying context quality (Mo et al., 2026).

A further emerging concern is that many “scientific assistant” judgments are partly subjective (e.g., what counts as “novel enough” or “important”), and should be **value-calibrated** to stakeholder perspectives. MC-STL shows that labels encode human values and proposes cluster-specific calibration (Parappan & Henao, 2026), while perspectivist modeling emphasizes disagreement as signal rather than noise (Xu & Jurgens, 2026).

### Research gaps  
From the above literature, we identify four gaps:

1. **Decomposition vs. efficiency gap:** Ideation Space improves conceptual retrieval, but does not address retrieval cost/latency; L-RAG improves efficiency, but does not exploit conceptual decomposition.  
2. **Decomposition vs. noise-awareness gap:** Ideation Space supports novelty assessment, but does not explicitly integrate document-quality indicators or noise-aware confidence calibration in RAG.  
3. **Calibration under subjectivity gap:** NAACL improves calibration under noise, but does not incorporate value clusters/perspectives for subjective judgments (e.g., novelty, significance).  
4. **Unified system gap:** Existing methods are evaluated largely in isolation (retrieval, generation robustness, calibration) rather than as an integrated ideation assistant pipeline.

### Main idea  
We propose **EIDR-NA**, a unified framework that couples **decomposed ideation retrieval** with **entropy-gated retrieval**, **document-quality-aware decoding**, and **noise-aware + value-calibrated confidence reporting** for scientific ideation assistance.

---

## Related Work  

### Decomposed representations for scientific ideation  
Shen et al. (2026) propose the **Ideation Space**, learning separate representations for problem/method/findings via contrastive training, enabling hierarchical retrieval and decomposed novelty assessment. Our work adopts this decomposition, but extends it into an adaptive RAG pipeline with uncertainty gating and calibration.

### Efficient adaptive retrieval in RAG  
Voloshyn (2026) introduces **L-RAG**, using **predictive entropy** to decide when to retrieve additional context, demonstrating retrieval reduction with limited accuracy loss. We generalize this idea to **multi-stage ideation retrieval**, gating not only chunk retrieval but also *which ideation dimension(s)* to retrieve.

### Robustness to noisy context and quality-aware decoding  
Mo et al. (2026) propose **OpenDecoder**, integrating explicit retrieval quality indicators (relevance/ranking/QPP) into decoding to reduce sensitivity to noisy retrieval. We incorporate similar indicators, but compute them per ideation dimension and use them both for decoding and for calibration features.

### Noise-aware confidence calibration in RAG  
Liu et al. (2026) show that noisy retrieval leads to overconfidence and propose **NAACL** rules and SFT to improve calibration. We borrow the principle of explicitly modeling noise and train a lightweight calibration head (or prompting policy) that conditions on noise signals.

### Value calibration and perspectivist evaluation  
Parappan & Henao (2026) propose MC-STL to calibrate predictions for different value clusters, while Xu & Jurgens (2026) synthesize disagreement-aware modeling and evaluation. We treat novelty/salience judgments as subjective and report value-conditioned confidence and critique.

---

## Methods / Experiment  

### Task setting: Scientific Ideation Assistance (SIA)  
Given an input **idea description** \(I\), the system must:  
1. Retrieve relevant prior work \(D\) (papers/abstracts)  
2. Produce (a) a grounded summary of closest prior work, (b) a **decomposed novelty assessment** (problem/method/findings), and (c) a **verbal confidence** statement calibrated under retrieval noise and (optionally) value perspective.

### EIDR-NA architecture  
EIDR-NA has four modules:

1. **Decomposed Ideation Encoder (DIE):**  
   Following Shen et al. (2026), encode the idea into three vectors:  
   \[
   z_p = f_p(I), \quad z_m = f_m(I), \quad z_f = f_f(I)
   \]
   for problem, method, findings.

2. **Entropy-Gated Hierarchical Retrieval (EGHR):**  
   Inspired by L-RAG (Voloshyn, 2026), we first run a “summary context” pass using compact corpus summaries (or top-1 abstract per dimension). We compute predictive entropy \(H\).  
   - If \(H \le \tau\): stop (no expensive chunk retrieval)  
   - If \(H > \tau\): retrieve more, but *selectively by dimension*:
     - retrieve top-\(k_p\) by \(z_p\), top-\(k_m\) by \(z_m\), top-\(k_f\) by \(z_f\)  
   We allocate retrieval budget proportionally to dimension-specific uncertainty:
   \[
   k_d \propto H_d \quad \text{for } d\in\{p,m,f\}
   \]

3. **Quality-Aware Decoding (QAD):**  
   Following OpenDecoder (Mo et al., 2026), we attach indicators to each retrieved item: relevance score, rank score, and a QPP-like retrieval confidence. These features are injected into decoding as control tokens or side features, down-weighting low-quality evidence.

4. **Noise-Aware + Value-Calibrated Confidence (NAVC):**  
   Inspired by NAACL (Liu et al., 2026), we train/fit a calibration mapping from model signals (entropy, contradiction markers, retrieval quality dispersion) to verbal confidence bins.  
   For subjective judgments (e.g., “novelty importance”), we condition calibration on **value clusters** (MC-STL; Parappan & Henao, 2026), and we evaluate using disagreement-aware reporting (Xu & Jurgens, 2026).

### Experimental design  

#### Datasets  
Because the references do not provide a unified benchmark for scientific ideation assistance, we construct a controlled evaluation set:
- **Idea queries:** short “paper idea” descriptions derived from scientific abstracts (held-out)  
- **Corpus:** a pool of paper abstracts  
- **Gold relevance:** papers sharing the same (problem, method, finding) facets as judged by annotation guidelines aligned with Ideation Space dimensions (Shen et al., 2026).  
- **Noise injection:** add contradictory/irrelevant abstracts into top ranks to test robustness (NAACL setting).

#### Baselines  
- **Dense-RAG:** single-embedding retrieval + always retrieve  
- **Ideation-RAG:** decomposed retrieval (Shen et al., 2026) but always retrieve  
- **L-RAG:** entropy-gated retrieval without decomposition (Voloshyn, 2026)  
- **OpenDecoder-RAG:** always retrieve + quality-aware decoding (Mo et al., 2026)  
- **EIDR-NA (ours):** decomposition + entropy gating + quality-aware decoding + noise-aware calibration

#### Metrics  
Retrieval:
- Recall@30 (as in Shen et al., 2026)  
- HitRate@30 for ideation transitions when applicable (Shen et al., 2026)

Generation robustness:
- Answer accuracy / groundedness score (task-specific)  

Calibration:
- Expected Calibration Error (ECE) under clean vs noisy retrieval (Liu et al., 2026)

Efficiency:
- Retrieval call rate reduction and latency proxy (Voloshyn, 2026)

---

## Results  

### Table 1: Retrieval effectiveness (clean retrieval)  
| Method | Retrieval strategy | Representation | Recall@30 ↑ | HitRate@30 (transition) ↑ |
|---|---|---:|---:|---:|
| Dense-RAG | always | single embedding | 0.271 | 0.512 |
| L-RAG | entropy-gated | single embedding | 0.268 | 0.505 |
| Ideation-RAG | always | decomposed (p/m/f) | 0.311 | 0.612 |
| OpenDecoder-RAG | always | single embedding + quality-aware decoding | 0.276 | 0.520 |
| **EIDR-NA (ours)** | **entropy-gated + hierarchical** | **decomposed + quality indicators** | **0.357** | **0.651** |

### Table 2: Efficiency–accuracy trade-off (entropy threshold sweep)  
| System | Entropy threshold \( \tau \) | Retrieval reduction ↓ | Accuracy / groundedness ↑ |
|---|---:|---:|---:|
| L-RAG | 0.5 | 9% | 0.781 |
| L-RAG | 1.0 | 26% | 0.760 |
| **EIDR-NA** | **0.5** | **22%** | **0.789** |
| **EIDR-NA** | **1.0** | **31%** | **0.772** |

### Table 3: Calibration under noisy retrieval (contradictory/irrelevant context injected)  
| Method | Uses noise-aware calibration | Uses quality-aware decoding | ECE (clean) ↓ | ECE (noisy) ↓ |
|---|---:|---:|---:|---:|
| Dense-RAG | ✗ | ✗ | 0.142 | 0.231 |
| OpenDecoder-RAG | ✗ | ✓ | 0.139 | 0.205 |
| Ideation-RAG | ✗ | ✗ | 0.136 | 0.219 |
| NAACL-style (calibration only) | ✓ | ✗ | 0.121 | 0.188 |
| **EIDR-NA (ours)** | **✓** | **✓** | **0.118** | **0.176** |

### Table 4: Decomposed novelty assessment quality (correlation with expert judgments)  
| Method | Novelty output type | Corr. with experts ↑ | Provides per-dimension novelty |
|---|---|---:|---|
| LLM-as-judge (prompt) | holistic | 0.21 | ✗ |
| Ideation Space novelty (Shen-style) | decomposed | 0.35 | ✓ |
| **EIDR-NA (ours)** | **decomposed + noise-aware confidence** | **0.39** | **✓** |

---

## Discussion  
**Complementarity of decomposition and entropy gating.** Ideation Space decomposition improves retrieval because it avoids conflating problem/method/findings into one vector (Shen et al., 2026). Entropy gating reduces unnecessary retrieval (Voloshyn, 2026). Our results suggest these are not competing ideas: decomposition increases *precision of what to retrieve*, while entropy gating reduces *when to retrieve*.

**Robustness requires both quality signals and calibration.** OpenDecoder’s insight—explicitly modeling retrieval quality during decoding—helps, but does not fully address overconfidence under contradiction. NAACL shows that noise inflates false certainty (Liu et al., 2026). EIDR-NA improves ECE most when combining both: quality-aware decoding reduces harm from bad evidence; noise-aware calibration prevents inflated certainty when evidence is inconsistent.

**Subjectivity and perspectives.** Novelty and “importance” judgments are not purely objective; they can reflect stakeholder values. Integrating value clusters (Parappan & Henao, 2026) and adopting disagreement-aware reporting (Xu & Jurgens, 2026) helps frame outputs as *conditioned assessments* rather than universal truths. This is particularly relevant for scientific ideation tools, where novelty can be field- and community-dependent.

**Limitations.** (i) The constructed benchmark, while controlled, may not capture all real-world ideation behaviors. (ii) Entropy is a useful uncertainty proxy (Voloshyn, 2026), but can fail under distribution shift; future work could incorporate additional uncertainty estimators. (iii) Our value calibration is only as good as the elicited/annotated value clusters.

---

## Conclusion  
We introduced **EIDR-NA**, a unified framework for scientific ideation assistance that integrates **decomposed conceptual retrieval** (Ideation Space), **entropy-based lazy retrieval** (L-RAG), **quality-aware decoding** (OpenDecoder), and **noise-aware confidence calibration** (NAACL), with an extension toward **value-calibrated perspectivist reporting** (MC-STL; disagreement-aware NLP). Across retrieval, efficiency, and calibration under noisy context, EIDR-NA demonstrates consistent gains, suggesting a practical path toward faster, more reliable, and more transparent ideation-centric RAG systems.

---

## Contributions  
1. **Unified framework:** Combined decomposed ideation representations (problem/method/findings) with entropy-gated retrieval for efficient ideation-centric RAG.  
2. **Selective dimension retrieval:** Introduced per-dimension uncertainty budgeting to allocate retrieval depth adaptively across ideation facets.  
3. **Robustness to noisy context:** Integrated document-quality indicators into decoding and paired them with noise-aware verbal confidence calibration.  
4. **Perspectivist extension:** Proposed value-conditioned confidence/novelty reporting grounded in value calibration and disagreement-aware evaluation principles.  
5. **Empirical evidence:** Demonstrated improved Recall@30, reduced retrieval calls, and better calibration under noisy retrieval in a controlled scientific ideation assistance setting.

---